\zlabel{triptych:concise:themes:start}
\section{The Propers: Themes and Movement}
\runninghead{Themes and Movement}

On the First Sunday of Advent, the Roman Church opens her liturgical year at the Basilica of St.~Mary Major (\latin{Statio ad S.~Mariam maiorem}, ``Station at St.~Mary Major''), the Roman church of the Crib. Her ministers wear violet and keep silent the Christmas hymn \latin{Gloria in excelsis Deo} (``Glory to God in the highest''). Yet the appointed Gospel is not the Annunciation or Bethlehem, but Jesus' Passion-Week prophecy in the Temple: signs in sun, moon, and stars, nations distressed at the roaring sea, men withering in terror, and the Son of Man coming in a cloud with power and majesty (Luke~21:25--33). Beside it stand St.~Paul's call to put on the armour of light (Rom.~13:11--14a), a Sunday Alleluia begging mercy and salvation (Ps.~84:8), and a Communion antiphon promising goodness and fruitful earth (Ps.~84:13).

Three patristic and liturgical interpretations answer how cosmic upheaval, a dawn call to wake, a Marian basilica, and fruitful soil belong in one Mass. Together they match Christ's threefold coming (\latin{triplex adventus}), summarized by John Beleth (\work{Summa de ecclesiasticis officiis}~LXIV) and preached by St.~Bernard of Clairvaux (\work{Sermones in Adventu}~3.3--4, 5.1): in human flesh, invisibly in grace, and openly in majesty. The first reads the Mass before the Last Judgment, where an aging world's decay detaches hearts for bodily resurrection and cosmic renewal. The second reads it through the Incarnation and Mary, joining St.~Mary Major to the ``light cloud'' of Luke~21:27 and the virgin earth of Psalm~84:13. The third looks to the present hour between Bethlehem and Judgment, where prevenient grace, confession, moral ascent, and the Eucharist purify sinners for the coming feast. (For scriptural dates and settings, see \work{Scriptural Date and Location} above.)

\subsection{Entrance, Prayer, and the Pre-Dawn Summons}

The Introit opens with \latin{Ad te levavi animam meam} (``To thee have I lifted up my soul,'' Ps.~24:1--3, 4b [Heb.~25:1--3a, 4]). Against the Gallican Vulgate, it drops the heading and initial \latin{Domine} (``O Lord''), keeps Old Latin \latin{qui te exspectant} (``who wait on thee'') over \latin{qui sustinent te} (``who wait for thee''), and omits imprecatory verse~4a. Its verse sings \latin{Vias tuas, Domine, demonstra mihi} (``Show me thy ways, O Lord'') before \latin{Gloria Patri} (``Glory be to the Father'') and the repeated antiphon. John Beleth (\work{Summa}~LV) hears the Church looking from afar toward the Nativity. William Durandus (\work{Rationale divinorum officiorum}~VI.3) notes that the cantor raises the pitch from low to high to rouse sleepers. The psalm's opening cry draws three accents. St.~Robert Bellarmine hears a soul tearing itself from earthly ``briers and thorns'' to trust God alone. St.~Augustine hears the whole Christ (\latin{totus Christus}), Head speaking in His Church, praying against pride. St.~Thomas Aquinas, following St.~John Damascene, takes lifting the soul as prayer itself: \latin{ascensus intellectus in Deum} (``the ascent of the intellect into God'').

The Collect, \latin{Excita, quaesumus, Domine, potentiam tuam, et veni} (``Bestir, O Lord, Thy might, we pray Thee and come''), stands alone on this first-class Sunday, omitting St.~Saturninus on 29~November~2026 and the seasonal extra orations abolished in 1960. Its opening echoes Psalm~79:3 (Heb.~80:3), which Bellarmine expounds as begging God to rouse His hidden might. Where that psalm appeals against outward foes, the Collect names an inward peril: \latin{ab imminentibus peccatorum nostrorum periculis}, ``from approaching dangers brought on by our sins.'' Ending \latin{Qui vivis} (``Who livest''), it addresses Christ directly, attributing rescue to His protection (\latin{te protegente eripi}) and deliverance (\latin{te liberante salvari}). As Bl.~Ildefonso Schuster notes (\work{The Sacramentary}~I, p.~321), the Church begs our Lord to save her by His power from sin's dangers.

The Epistle (Rom.~13:11--14a) drops the Vulgate's opening connective to begin \latin{Fratres: Scientes quia hora est} (``Brethren: Knowing that it is now the hour''), and stops at \latin{sed induimini Dominum Jesum Christum} (``but put ye on the Lord Jesus Christ''), omitting verse~14b's warning against making provision for the flesh. St.~Paul bids believers wake, cast off works of darkness, put on the armour of light, and abandon revelry, drunkenness, impurity, and faction. St.~John Chrysostom preaches that call before Resurrection and Judgment, comparing Christians to runners straining near the finish. Durandus links putting on Christ to the Incarnation, because the Son became man so man might become a son of God. In \work{Confessions}~VIII.12.29, the passage broke St.~Augustine's divided will, though his conversion turned on reading through omitted verse~14b.

\subsection{Waiting, Mercy, and the Coming Son of Man}

After the Epistle, the Gradual (Ps.~24:3b, 4b [Heb.~25:3a, 4]) returns to Psalm~24, dropping \latin{etenim} (``for''), adding \latin{Domine}, and keeping Old Latin \latin{notas fac mihi} (``make known to me'') beside the Introit verse's \latin{demonstra mihi} (``show me''). Both verses ask separately for God's ``ways'' (\latin{vias}) and ``paths'' (\latin{semitas}), which Aquinas distinguishes as the highway of commandments and the narrower footpath of evangelical counsels.

Next, the Sunday Mass sings \latin{Alleluia} with Psalm~84:8 (Heb.~85:8; modern English~85:7), \latin{Ostende nobis, Domine, misericordiam tuam: et salutare tuum da nobis} (``Shew us, O Lord, thy mercy; and grant us thy salvation''), omitted on resumed Advent weekdays. Why does \latin{Alleluia} sound while \latin{Gloria in excelsis} stays silent? Durandus (\work{Rationale}~VI.2) answers that Advent is ``partly of rejoicing and partly of mourning'' (\latin{partim exultationis, et partim moeroris}), withholding the angels' hymn so it may be sung more devoutly on Christmas night. Beleth (\work{Summa}~LXIV) calls the season ``partly of gladness'' (\latin{partim laetitiae}), explaining that \latin{Gloria} (and, in medieval usage, \latin{Ite, missa est}) is omitted because the Nativity and Resurrection are eagerly awaited. On Psalm~84:8, Bellarmine hears \latin{salutare tuum} directly of the Incarnate Son. Augustine stresses the verb \emph{show}: God reveals His mercy to cure human boasting so man recognizes every good as His gift.

Two reversals stand in the Gospel (Luke~21:25--33). After men wither in fear, Jesus bids disciples lift up their heads (\latin{levate capita vestra}, v.~28) because redemption is near. Where Matthew~24:32 names only the fig tree, Luke adds ``and all the trees'' (\latin{et omnes arbores}, v.~29), before Christ pledges that His words will outlast heaven and earth (v.~33). St.~Gregory the Great reads the signs literally as heralds of the strict Judge and takes the fig tree's fruit ironically as the world's ruin. St.~Cyril of Alexandria anchors redemption in bodily resurrection. On verse~33, St.~Bede explains cosmic renewal in outward form (\latin{per imaginem}), drawing on St.~Augustine (\work{De civitate Dei}~XX). St.~Ambrose reads the signs symbolically and identifies the cloud of verse~27 with the prophets and Mary. For him, the fig tree and all the trees are Israel and every nation blossoming in faith, while Christ also comes within each soul (\latin{in singulis}). The Sunday \latin{Credo} follows, confessing both the Incarnation of Mary and the returning Judge.

\subsection{Altar, Harvest, and Preparation for the Feast}

At the Offertory (Ps.~24:1--3 [Heb.~25:1--3a]), the choir repeats the Introit antiphon verbatim, which Rupert of Deutz (\work{De divinis officiis}~III.2) reads as the Church's continuing cry toward the heavenly Jerusalem. The Secret (\latin{Haec sacra nos, Domine}) begs that these Holy Mysteries cleanse worshippers by their mighty power (\latin{potenti virtute mundatos}) to approach their divine Source (\latin{principium}) more purely. Schuster (\work{The Sacramentary}~I, p.~322) hears in it the ancient Roman epiklesis. Without a proper Advent Preface, Sundays take the Preface of the Most Holy Trinity and resumed weekdays take the Common Preface. The Trinity Preface confesses distinctness in Persons, oneness in essence, and equality in majesty beside angelic praise.

At Communion, the choir sings Psalm~84:13 (Heb.~85:13; modern English~85:12) without \latin{Etenim} (``For''): \latin{Dominus dabit benignitatem: et terra nostra dabit fructum suum} (``The Lord will give goodness: and our earth shall yield her fruit''). St.~Robert Bellarmine and Sicard of Cremona identify that goodness with the Holy Spirit, through whom the Virgin Mary yields Christ. St.~Augustine, reading Old Latin \latin{suavitatem} (``sweetness''), expounds the verse instead of sinners' hearts ploughed by confession. Bl.~Ildefonso Schuster applies it to Holy Communion refreshing arid souls. The Postcommunion (\latin{Suscipiamus, Domine}) turns Psalm~47:10's past tense (Heb.~48:10; modern English~48:9) into a plea for temple mercy to prepare with fitting honour for ``the coming solemnities of our reparation'' (\latin{reparationis nostrae ventura solemnia}), before \latin{Ite, missa est} (``Go, the Mass is ended'') and the Last Gospel (John~1:1--14).
\zlabel{triptych:concise:themes:end}
